% =============================================================================
%  ملخّص الدرس — الأثر الكتابي على الكراس (الدرس 01)
%  المجال 1: بيئة التعامل مع الحاسوب — الوحدة 1: تقنية المعلومات (IT)
%  الموضوع: المفاهيم القاعدية لتقنية المعلومات
%  السنة الدراسية: 2026 / 2027
% =============================================================================
%  البناء (من هذا المجلد):
%    xelatex -interaction=nonstopmode -halt-on-error -output-directory=build lesson-summary.tex
% =============================================================================

\documentclass[11pt, a4paper]{article}

% ── المشترك (التمهيد، الألوان، الترويسة، التذييل، الأدوات) ───────────────────
\input{"../../shared/common.tex"}

% ── خاص بهذا المستند: عنوان الملخّص ─────────────────────────────────────────
\newcommand{\summarytitle}[3]{%
  \begin{center}
    {\large\bfseries\color{primary} ملخّص الدرس — الأثر الكتابي على الكراس}
  \end{center}
  \vspace{1pt}
  {\color{secondary}\hrule height 1.2pt}
  \vspace{4pt}
  {\cardsetup
  \noindent
  \begin{tabularx}{\textwidth}{@{}>{\bfseries\color{primary}}p{3.2cm}X@{}}
    \toprule
    المجال التعلمي & #1 \\
    الوحدة التعلمية & #2 \\
    الموضوع & #3 \\
    \bottomrule
  \end{tabularx}}\vspace{10pt}
}

% ── خاص بهذا المستند: محور ──────────────────────────────────────────────────
\newcommand{\summaryaxis}[2]{%
  \par\vspace{6pt}%
  {\normalsize\bfseries\color{primary} #1. #2}\par
  \vspace{-0.3em}{\color{secondary}\hrule height 0.8pt}\par\vspace{3pt}%
}

% ── خاص بهذا المستند: مصطلح ────────────────────────────────────────────────
\newcommand{\term}[2]{%
  \textbullet\ \textbf{\color{primary}#1:} #2\par\vspace{2pt}%
}

% ── خاص بهذا المستند: صندوق الخلاصة ────────────────────────────────────────
\newtcolorbox{leadbox}[1][]{
  enhanced, colback=lightbg, colframe=primary,
  coltitle=white, fonttitle=\bfseries, title=#1,
  boxrule=1pt, arc=4pt, left=10pt, right=10pt, top=8pt, bottom=8pt,
  breakable,
}

% ══════════════════════════════════════════════════════════════════════════════
\begin{document}

\officialheader

\summarytitle{بيئة التعامل مع الحاسوب}{تقنية المعلومات (IT)}{المفاهيم القاعدية لتقنية المعلومات}

% ══════════════════════════════════════════════════════════════════════════════
\summaryaxis{1}{المفاهيم الأولية}

\term{التقنية (Technique)}{هي المهارة والقدرة الفردية على ممارسة الأعمال والأنشطة لإنجاز منتج معيّن.}

\term{التكنولوجيا (Technologie)}{كلمة ذات أصل يوناني مركّبة من مقطعين: (Techno: مهارة فنية) + (Logos: علم)، وتعني علم المهارات والتقنيات التطبيقية والاستغلال المنطقي للمعارف العلمية.}

\term{الإعلام (Information)}{هو تقديم الأخبار والمعلومات ونقل رسالة من مرسل إلى مستقبل.}

\term{الاتصال (Communication)}{هو عملية التفاعل ونقل الأفكار والمعلومات بين مرسل ومستقبل في سياق معيّن.}

% ══════════════════════════════════════════════════════════════════════════════
\summaryaxis{2}{مفهوم تكنولوجيا الإعلام والاتصال (TIC)}

\noindent هي التلاقي والتزاوج بين عتاد الحواسيب، البرمجيات، وشبكات الاتصال لجمع البيانات وتخزينها ومعالجتها ونقلها.

\vspace{4pt}
\noindent\textbf{\color{secondary}مكوّناتها الأساسية:}
\begin{itemize}[label=\textbullet, leftmargin=1.2em, itemsep=2pt, topsep=2pt, parsep=0pt]
  \item \textbf{تقنيات المعالجة والتخزين:} الأجهزة (كمبيوتر، هواتف) والبرامج (أنظمة تشغيل، تطبيقات).
  \item \textbf{شبكات الاتصال:} شبكات الاتصال السلكية واللاسلكية والسمعي البصري.
\end{itemize}

% ══════════════════════════════════════════════════════════════════════════════
\summaryaxis{3}{البيانات والمعلومات والمعالجة}

\term{البيانات (Data / Données)}{هي الحقائق والأرقام والرموز الأولية الخام التي لا تعطي معنى واضحًا بمفردها.}

\term{المعلومات (Informations)}{هي البيانات التي جرت معالجتها وتنظيمها فأصبحت ذات معنى وسياق مفيد.}

\term{المعالجة (Processing / Traitement)}{هي مجموعة العمليات الحسابية والمنطقية المطبّقة على البيانات لتحويلها إلى معلومات.}

\vspace{4pt}
\begin{center}
\begin{tikzpicture}[
  node distance=2.8cm,
  box/.style={draw=primary, fill=lightbg, thick, rounded corners=3pt,
              minimum width=3.4cm, minimum height=1.0cm, align=center},
  ar/.style={-{Latex[length=3mm]}, thick, primary}
]
  \node[box] (d) {بيانات\\\footnotesize (Data)};
  \node[box, right=of d] (p) {معالجة\\\footnotesize (Processing)};
  \node[box, right=of p] (i) {معلومات\\\footnotesize (Information)};
  \draw[ar] (d) -- (p);
  \draw[ar] (p) -- (i);
\end{tikzpicture}
\end{center}

% ══════════════════════════════════════════════════════════════════════════════
\summaryaxis{4}{مفهوم المعلوماتية (Informatique)}

\term{أصل الكلمة}{كلمة فرنسية مركّبة من جزأين: \textbf{INFORMATION} (معلومات) + \textbf{AUTOMATIQUE} (آلية).}

\term{التعريف}{هي علم المعالجة الآلية للمعلومات باستخدام جهاز الحاسوب والبرمجيات المخصّصة لذلك.}

\vspace{6pt}
\begin{leadbox}[الخلاصة المرمزة]
  \begin{center}
    INFORMATION (معلومات) $+$ AUTOMATIQUE (آلية) $\;=\;$ \textbf{\color{primary}INFORMATIQUE (المعلوماتية)}
  \end{center}
\end{leadbox}

\end{document}
